\section*{Chapitre \ref{ch:b2:ligand-field} --- Théorie du champ des ligands et couleur}

\begin{solution}{exo:b2:ligand-field:1}
$d^3$~: $t_{2g}^3$, ESCC $-1.2\Delta_o$. $d^8$~: $t_{2g}^6e_g^2$, ESCC $-2.4 + 1.2 = -1.2\Delta_o$.
\end{solution}

\begin{solution}{exo:b2:ligand-field:2}
\ce{[Fe(H2O)6]^2+}, $d^6$ \omterm{def:b2:ligand-field:spin}{haut spin} $t_{2g}^4e_g^2$~: quatre. \ce{[Fe(CN)6]^4-}, \omterm{def:b2:ligand-field:spin}{bas spin}
$t_{2g}^6$~: aucun (\omterm{def:b2:diatomic-mos:magnetism}{diamagnétique}).
\end{solution}

\begin{solution}{exo:b2:ligand-field:3}
600 nm correspond à la lumière orange~; le complexe paraît bleu (bleu verdâtre).
\end{solution}

\begin{solution}{exo:b2:ligand-field:4}
$\ce{Cl-} < \ce{H2O} < \ce{NH3} < \ce{CN-}$.
\end{solution}

\begin{solution}{exo:b2:ligand-field:5}
Eau~: $\Delta_o < P$, \omterm{def:b2:ligand-field:spin}{haut spin}, $t_{2g}^3e_g^2$, cinq électrons célibataires. Cyanure~: $\Delta_o > P$,
\omterm{def:b2:ligand-field:spin}{bas spin}, $t_{2g}^5$, un électron célibataire.
\end{solution}

\begin{solution}{exo:b2:ligand-field:6}
Dans le complexe tétraédrique \ce{[CoCl4]^2-}, l'éclatement vaut environ les $4/9$ d'un éclatement octaédrique, et
le chlorure est un ligand plus faible que l'eau~: la bande d--d se déplace vers les grandes longueurs d'onde (de l'orange
au rouge), et le complexe paraît bleu. Les complexes tétraédriques absorbent aussi plus fortement,
d'où la couleur intense.
\end{solution}

\begin{solution}{exo:b2:ligand-field:7}
\ce{Ni^2+} est $d^8$. Plan carré~: les quatre orbitales inférieures contiennent les huit électrons appariés,
$\mathrm d_{x^2-y^2}$ reste vide~: \omterm{def:b2:diatomic-mos:magnetism}{diamagnétique}. Tétraédrique~: $e^4t_2^4$, deux électrons célibataires
dans l'ensemble $t_2$.
\end{solution}

\begin{solution}{exo:b2:ligand-field:8}
$\tilde\nu = 1/(500 \times 10^{-7}\,\text{cm}) = \qty{20000}{cm^{-1}}$~; $\Delta_o =
0.1196/(500 \times 10^{-9}) = \qty{2.39e5}{J/mol}$, \qty{239}{kJ/mol}.
\end{solution}

\begin{solution}{exo:b2:ligand-field:9}
\ce{Zn^2+} est $d^{10}$ et \ce{Sc^3+} $d^0$~: aucune \omterm{def:b2:ligand-field:d-d}{transition d--d} n'est possible (pas de niveau d vide
pour \ce{Zn^2+}, pas d'électron d pour \ce{Sc^3+}), donc pas d'absorption dans le visible.
\end{solution}

\begin{solution}{exo:b2:ligand-field:10}
Co(III) $d^6$ + six doublets de l'ammoniac~: 18 électrons. Douze remplissent les six \omterm{def:b2:diatomic-mos:bonding}{orbitales liantes},
six remplissent les \omterm{def:b2:diatomic-mos:bonding}{orbitales non liantes} $t_{2g}$ (l'ammoniac crée un champ assez fort pour un \omterm{def:b2:ligand-field:spin}{bas spin} avec
le cobalt(III))~; $e_g^*$ est vide. Tous les électrons sont appariés~: \omterm{def:b2:diatomic-mos:magnetism}{diamagnétique}.
\end{solution}

\begin{solution}{exo:b2:ligand-field:11}
\ce{CO} est un bon donneur $\sigma$ et surtout un accepteur $\pi$~: ses orbitales $\pi^*$ vides
abaissent le niveau $t_{2g}$ par \omterm{def:b2:ligand-field:pi-ligands}{rétrodonation}, ce qui élargit $\Delta_o$. \ce{F-} est un donneur $\pi$, qui
élève $t_{2g}$ et réduit $\Delta_o$. Le modèle des charges ponctuelles ne voit que la charge et se trompe
sur l'ordre.
\end{solution}

\begin{solution}{exo:b2:ligand-field:12}
Sans champ, les enthalpies d'hydratation suivraient une courbe régulière (ions de plus en plus petits le long de
la série). L'eau donne des aquacomplexes \omterm{def:b2:ligand-field:spin}{haut spin} dont l'ESCC est nulle pour $d^0$, $d^5$, $d^{10}$
et maximale pour $d^3$ et $d^8$~: cette stabilisation supplémentaire ajoute deux bosses, culminant vers
\ce{V^2+} et \ce{Ni^2+}, au-dessus de la droite passant par \ce{Ca^2+}, \ce{Mn^2+}, \ce{Zn^2+}.
\end{solution}

\begin{solution}{pb:b2:ligand-field:1}
\textbf{1.} $d^3$.
\textbf{2.} $t_{2g}^3$, un électron dans chacune des trois orbitales, spins parallèles.
\textbf{3.} Non~: trois électrons tiennent dans $t_{2g}$ sans appariement.
\textbf{4.} $-1.2\Delta_o$.
\textbf{5.} Trois.
\textbf{6.} L'ESCC est grande et les orbitales $e_g^*$, qui pointent vers les ligands, sont
vides~: retirer ou ajouter un ligand coûte une grande partie de cette stabilisation.
\textbf{7.} Six ions oxyde aux sommets d'un octaèdre légèrement déformé.
\textbf{8.} $1/(556 \times 10^{-7}) = \qty{17990}{cm^{-1}}$, environ \qty{1.80e4}{cm^{-1}}.
\textbf{9.} $0.1196/(556 \times 10^{-9}) = \qty{2.15e5}{J/mol}$~: \qty{215}{kJ/mol}.
\textbf{10.} Le jaune-vert (556 nm) et le violet (405 nm).
\textbf{11.} Du rouge, avec un peu de bleu~: le rouge profond du rubis.
\textbf{12.} Un ion $d^3$ possède plusieurs configurations excitées avec un électron dans $e_g$~;
la répulsion entre électrons leur donne des énergies différentes, d'où plusieurs bandes (traité dans le
volume de troisième année).
\textbf{13.} $0.1196/(620 \times 10^{-9}) = \qty{1.93e5}{J/mol}$, \qty{193}{kJ/mol}
(\qty{16130}{cm^{-1}}).
\textbf{14.} Le rubis~: son éclatement est plus grand.
\textbf{15.} Vers l'orange-rouge~: c'est maintenant le rouge qui est absorbé, et le vert (avec un peu de bleu) est transmis.
\textbf{16.} Dans le béryl, les sites du chrome sont un peu plus grands et les ions oxyde plus éloignés,
et l'environnement diffère~: le champ est plus faible.
\textbf{17.} Non~: un ion $d^3$ garde trois électrons célibataires dans tout champ octaédrique.
\textbf{18.} $d^9$.
\textbf{19.} L'aquacomplexe absorbe dans le rouge et l'orange par une \omterm{def:b2:ligand-field:d-d}{transition d--d}~: la
lumière transmise est bleue.
\textbf{20.} Il devient blanc~: sans les ligands aqua autour du cuivre (le sulfate prend leur
place), le champ est plus faible et la bande se déplace dans l'infrarouge.
\textbf{21.} L'eau se lie de nouveau au cuivre et l'aquacomplexe bleu se reforme.
\textbf{22.} À plus courte longueur d'onde~: l'ammoniac, au-dessus de l'eau dans la série, donne un plus grand
éclatement (le bleu-violet profond du complexe ammine).
\textbf{23.} Non~: $t_{2g}^6e_g^3$ est la seule configuration.
\textbf{24.} Oui~: un électron célibataire.
\textbf{25.} $\boldsymbol{\Delta_o \approx \qty{1.8e4}{cm^{-1}}}$, environ
$\boldsymbol{\qty{215}{kJ/mol}}$.
\end{solution}
